% Reporte Hormgas de Langton
\documentclass[12pt,a4paper]{article}
\usepackage[utf8]{inputenc}
\usepackage[T1]{fontenc}
\usepackage{graphicx}
\usepackage{xcolor}
\usepackage{listings}
\usepackage[strings]{underscore}
\usepackage{hyperref}
\usepackage{geometry}
\geometry{margin=2.5cm}
\lstset{
  language=Python,
  basicstyle=\ttfamily\small,
  breaklines=true,
  breakatwhitespace=true,
  columns=fullflexible,
  showstringspaces=false,
  keywordstyle=\bfseries,
  commentstyle=\itshape,
  stringstyle=\color{blue}
}

\title{Reporte Hormigas de Langton}
\author{Eduardo Hernández Castro}
\date{Complex Sytems -- GENARO JUAREZ MARTINEZ -- Grupo 7CM4}

\begin{document}
\maketitle

\section{Introducción}
En este documento se registran las observaciones y fundamentos técnicos de la simulación de hormigas basada en las reglas de Langton y el comportamiento de los distintos tipos de hormigas.

\section{Fundamento del programa}
La simulación está diseñada para evaluar cómo interactúan cuatro tipos de hormigas en una grilla con reglas de movimiento de Langton y ciclo de vida controlado. Cada sección del código aporta una parte clave del modelo.

\subsection{Reglas principales}
\begin{itemize}
  \item Cada hormiga vive hasta 80 iteraciones. Esta regla se define en \texttt{config.py} mediante la constante \texttt{LIFESPAN} y se aplica en \texttt{simulation.py} durante la fase de envejecimiento y muerte.
  \item El estado de cada celda es negro/blanco, pero ese estado se maneja internamente y no se muestra en la visualización. El estado oculto se representa en \texttt{grid.py} con el arreglo \texttt{cell_state}.
  \item En celda negra la hormiga gira a la izquierda; en celda blanca gira a la derecha. Esta lógica está en \texttt{grid.py} dentro de la función \texttt{apply_langton_rule}.
  \item Después de girar, la hormiga invierte el color de la celda y avanza una celda. El volteo de la celda se hace con \texttt{flip_cell_state} en \texttt{grid.py} y el movimiento con \texttt{move_ant} o \texttt{attempt_movement_with_collision}.
  \item Si la celda destino está ocupada, la hormiga prueba hasta tres direcciones alternativas de forma aleatoria. Esta reintención se gestiona en \texttt{grid.py} dentro de \texttt{attempt_movement_with_collision}.
  \item Si todas las direcciones libres están ocupadas, la hormiga espera y no se mueve en esa iteración. El ant permanece en su posición cuando \texttt{attempt_movement_with_collision} no encuentra destino libre.
  \item Una hormiga nueva nace cuando una reproductora y una reina se enfrentan en direcciones opuestas y se cumple la probabilidad de reproducción. Esta interacción se evalúa en \texttt{simulation.py} durante la fase de reproducción, usando \texttt{REPRODUCTION_PROBABILITY} y \texttt{OPPOSITE_DIRECTIONS} desde \texttt{config.py}.
  \item Se detectan escenarios de aislamiento, sobrepoblación y estabilidad con umbrales de población y ocupación. Estas condiciones se definen en \texttt{config.py} y se usan en \texttt{simulation.py} dentro de \texttt{detect_scenario}.
\end{itemize}

\section{Estructura del código}
La aplicación está organizada en módulos claros que separan la lógica de configuración, el modelo de datos, la simulación y la visualización.

\subsection{Configuración}
El archivo \texttt{config.py} contiene las constantes que determinan el comportamiento general:
\begin{itemize}
  \item \texttt{GRID_SIZE}, \texttt{OCCUPANCY_RATIO} y \texttt{MAX_ANTS} para la dimensión y el número de hormigas iniciales.
  \item \texttt{ANT_DENSITIES} para la distribución inicial de los tipos de hormiga: 1\% reinas, 55\% trabajadoras, 9\% reproductoras y 35\% soldados.
  \item \texttt{LIFESPAN} define las 80 iteraciones de vida máxima de cada hormiga.
  \item \texttt{REPRODUCTION_PROBABILITY} y \texttt{OPPOSITE_DIRECTIONS} controlan cuándo se puede generar una nueva hormiga.
  \item \texttt{ISOLATION_THRESHOLD}, \texttt{OVERPOPULATION_GRID_THRESHOLD}, \texttt{OVERPOPULATION_DEATH_RATE}, \texttt{STABLE_POPULATION_VARIANCE} y \texttt{STABLE_GENERATIONS_REQUIRED} fijan los umbrales de escenario.
\end{itemize}

\subsection{Modelo de hormigas}
El módulo \texttt{ant.py} define la clase base \texttt{Ant} y las subclases \texttt{Queen}, \texttt{Worker}, \texttt{Reproducer} y \texttt{Soldier}. Cada hormiga mantiene su posición, dirección, edad y color de visualización. El método \texttt{age_one_iteration()} incrementa la edad y \texttt{is_alive()} determina si la hormiga debe morir cuando supera la longevidad.

\subsection{Grilla y reglas de Langton}
El módulo \texttt{grid.py} implementa la grilla de ocupación y el estado oculto de las celdas. Los componentes más relevantes son:
\begin{itemize}
  \item \texttt{occupancy}: matriz booleana que indica qué celdas tienen hormigas.
  \item \texttt{cell_state}: matriz de valores 0/1 que representa el color oculto de cada celda.
  \item \texttt{apply_langton_rule}: aplica la regla de giro y calcula la nueva posición.
  \item \texttt{flip_cell_state}: cambia el color de la celda actual después del giro.
  \item \texttt{move_ant}: mueve la hormiga a una celda vacía.
  \item \texttt{attempt_movement_with_collision}: gestiona choques y reintentos en hasta tres direcciones adicionales.
  \item \texttt{get_occupancy_ratio}: calcula el porcentaje de celdas ocupadas por hormigas.
\end{itemize}

\subsection{Simulación}
El módulo \texttt{simulation.py} coordina la ejecución completa:
\begin{itemize}
  \item \texttt{_initialize_ants}: crea la configuración inicial, asignando posiciones aleatorias y tipos de hormiga según \texttt{ANT_DENSITIES}.
  \item \texttt{iterate}: ejecuta el ciclo de movimiento, reproducción, envejecimiento y muerte por generación.
  \item \texttt{detect_scenario}: clasifica el estado actual como aislamiento, sobrepoblación o estabilidad.
  \item \texttt{run}: itera el número especificado de generaciones.
  \item Estadísticas por generación: total de hormigas, conteos por tipo, nacimientos, muertes, edad promedio y ocupación.
\end{itemize}

\subsection{Visualización y reporte}
El módulo \texttt{visualization.py} produce la salida gráfica y los resúmenes:
\begin{itemize}
  \item Gráficos de evolución poblacional y ocupación.
  \item Guardado de imágenes iniciales y finales del grid.
  \item Exportación de un archivo de resumen \texttt{summary.txt} con las métricas clave de cada corrida.
\end{itemize}

\subsection{Entrada principal}
El archivo \texttt{main.py} expone el CLI para ejecutar la simulación con parámetros ajustables:
\begin{itemize}
  \item número de iteraciones,
  \item tamaño de la grilla,
  \item ocupación inicial,
  \item modo toroidal,
  \item generación de reportes y CSV,
  \item registro de salida en archivo.
\end{itemize}

\section{Escenarios}
Según el código en \texttt{simulation.py} y los umbrales definidos en \texttt{config.py}, los escenarios se clasifican así:
\begin{enumerate}
  \item Aislamiento: ocurre cuando la población actual baja por debajo del 10\% de la población inicial. En nuestro caso, con 5000 hormigas iniciales, aislamiento se detecta si quedan menos de 500 hormigas.
  \item Sobrepoblación: ocurre cuando la ocupación de la grilla excede 80\% y la tasa de muerte de la iteración es mayor al 15\%. Esto implica una saturación del espacio combinada con un aumento rápido de defunciones.
  \item Sistema estable: ocurre cuando el total de hormigas se mantiene con una varianza relativa de 5\% o menos durante al menos 50 generaciones consecutivas.
\end{enumerate}

\section{Resultados}

\subsection{Caso 1 -- Seed 29976446, 100x100, 200 generaciones}
\begin{itemize}
  \item Archivo de resumen: \texttt{S29976446_G100_I200_O0.5_toroidal_20260605_074107_summary.txt}
  \item Generaciones: 200
  \item Hormigas finales: 1 (soldado)
  \item Ocupación final: 0.01\%
  \item Hormigas en generación 80: 147
  \item Escenario detectado: aislamiento
\end{itemize}
\begin{figure}[h!]
  \centering
  \begin{minipage}[b]{0.48\textwidth}
    \centering
    \includegraphics[width=\textwidth]{output/S29976446_G100_I200_O0.5_toroidal_20260605_074107_initial.png}
    \caption*{Caso 1: estado inicial}
  \end{minipage}
  \hfill
  \begin{minipage}[b]{0.48\textwidth}
    \centering
    \includegraphics[width=\textwidth]{output/S29976446_G100_I200_O0.5_toroidal_20260605_074107_final.png}
    \caption*{Caso 1: estado final}
  \end{minipage}
\end{figure}

\subsection{Caso 2 -- Seed 401529739, 100x100, 250 generaciones}
\begin{itemize}
  \item Archivo de resumen: \texttt{S401529739_G100_I250_O0.5_toroidal_20260605_075905_summary.txt}
  \item Generaciones: 250
  \item Hormigas finales: 1 (trabajadora)
  \item Ocupación final: 0.01\%
  \item Hormigas en generación 80: 183
  \item Escenario detectado: aislamiento
\end{itemize}
\begin{figure}[h!]
  \centering
  \begin{minipage}[b]{0.48\textwidth}
    \centering
    \includegraphics[width=\textwidth]{output/S401529739_G100_I250_O0.5_toroidal_20260605_075905_initial.png}
    \caption*{Caso 2: estado inicial}
  \end{minipage}
  \hfill
  \begin{minipage}[b]{0.48\textwidth}
    \centering
    \includegraphics[width=\textwidth]{output/S401529739_G100_I250_O0.5_toroidal_20260605_075905_final.png}
    \caption*{Caso 2: estado final}
  \end{minipage}
\end{figure}

\subsection{Notas de terminal y logs}
El log \texttt{output/logs/find_seed_20260605_074103.txt} confirma la búsqueda de semilla y los resultados:
\begin{itemize}
  \item Grid inicial 100x100 con 5000 hormigas y 50\% de ocupación.
  \item Al finalizar 200 iteraciones quedó 1 hormiga y 0.01\% de ocupación.
  \item En iteración 80 se tenía 147 hormigas.
\end{itemize}

\subsection{Dependencia de la semilla y el tamaño del grid}
La configuración inicial depende de la semilla aleatoria y del tamaño del grid. Al ser la grilla toroidal, dos tamaños de grid diferentes no pueden producir la misma configuración inicial válida, ya que la estructura de vecinos y la cantidad de celdas difieren.

\section{Conclusiones}
La evolución del sistema depende fuertemente de la semilla pseudoaleatoria usada para generar el estado inicial y de la incertidumbre inherente a respetar manualmente la ocupación y las proporciones de cada tipo de hormiga. En un grid de 100x100 con 50\% de ocupación, colocar a mano miles de hormigas manteniendo las proporciones de reinas, trabajadoras, reproductoras y soldados no es práctico, y por eso el estado inicial se construye aleatoriamente a partir de una semilla. Esto hace muy difícil encontrar un estado inicial ideal en el que, por ejemplo, varias reproductoras queden correctamente ubicadas alrededor de las reinas y puedan reproducirse cada pocas iteraciones.

En los casos estudiados, incluso cuando una semilla produce más hormigas en la iteración 80, la ubicación de las hormigas que sobreviven más allá de las originales resulta decisiva. Las posiciones relativas de estas hormigas determinan si seguirán encontrándose con reproductoras y reinas, y por eso la diferencia entre los resultados de 200 iteraciones y 250 iteraciones no solo está en cuántas hormigas quedan, sino en cómo quedan distribuidas. Esto se aprecia en los archivos de salida \texttt{output/S401529739_G100_I250_O0.5_toroidal_20260605_075905.png} y \texttt{output/S29976446_G100_I200_O0.5_toroidal_20260605_074107.png}.

\begin{figure}[h!]
  \centering
  \includegraphics[width=0.8\textwidth]{output/S401529739_G100_I250_O0.5_toroidal_20260605_075905.png}
  \caption*{Caso 2: estado final en 250 generaciones}
\end{figure}
\begin{figure}[h!]
  \centering
  \includegraphics[width=0.8\textwidth]{output/S29976446_G100_I200_O0.5_toroidal_20260605_074107.png}
  \caption*{Caso 1: estado final en 200 generaciones}
\end{figure}

Otro efecto importante es que, en la mayoría de los casos, las hormigas nacidas después de la primera generación quedan demasiado separadas entre sí. Ese gran distanciamiento complica mucho que vuelvan a encontrarse más adelante, porque el movimiento de Langton tiende a dispersarlas y no a atraerlas. Por lo tanto, aunque haya un repunte en las generaciones intermedias, la distancia entre nacimientos reduce la probabilidad de que ocurran nuevos encuentros reproductivos y favorece la deriva hacia el aislamiento.

Estas observaciones refuerzan que el sistema es sensible a la configuración inicial y que los estados verdaderamente estables o de sobrepoblación son difíciles de lograr con generación aleatoria simple. Para avanzar, sería necesario explorar estrategias que controlen mejor la ubicación inicial de los individuos clave (reinas y reproductoras) o que introduzcan mecanismos adicionales para mantener las hormigas en áreas de encuentro más frecuentes.

\section{Anexos}
En esta sección se incluyen los códigos completos de los módulos principales de la simulación.

\subsection{Archivo \texttt{config.py}}
\lstinputlisting[language=Python,caption={config.py}]{config.py}

\subsection{Archivo \texttt{ant.py}}
\lstinputlisting[language=Python,caption={ant.py}]{ant.py}

\subsection{Archivo \texttt{grid.py}}
\lstinputlisting[language=Python,caption={grid.py}]{grid.py}

\subsection{Archivo \texttt{simulation.py}}
\lstinputlisting[language=Python,caption={simulation.py}]{simulation.py}

\subsection{Archivo \texttt{main.py}}
\lstinputlisting[language=Python,caption={main.py}]{main.py}

\subsection{Archivo \texttt{visualization.py}}
\lstinputlisting[language=Python,caption={visualization.py}]{visualization.py}

\end{document}
